\documentclass[11pt]{article}

\usepackage[T1]{fontenc}
\usepackage{amsmath,amssymb,amsthm}
\usepackage{booktabs}
\usepackage{geometry}
\usepackage{hyperref}
\usepackage{xurl}

\geometry{margin=1in}
\hypersetup{
  colorlinks=true,
  linkcolor=blue,
  citecolor=blue,
  urlcolor=blue,
  pdftitle={Excluding Parameter Three at Order Thirteen in the gamma--theta Conjecture},
  pdfauthor={Alec Kriebel},
  pdfsubject={One-guard eternal domination and clique covering},
  pdfkeywords={eternal domination, domination number, clique cover, gamma--theta conjecture, SAT certificates}
}

\newtheorem{theorem}{Theorem}[section]
\newtheorem{lemma}[theorem]{Lemma}
\newtheorem{proposition}[theorem]{Proposition}
\newtheorem{corollary}[theorem]{Corollary}
\theoremstyle{definition}
\newtheorem{definition}[theorem]{Definition}
\theoremstyle{remark}
\newtheorem{remark}[theorem]{Remark}

\newcommand{\gaminf}{\gamma^{\infty}}
\newcommand{\V}{V}
\newcommand{\F}{\mathcal F}
\newcommand{\sha}{\operatorname{SHA256}}
\newcommand{\filehash}[1]{\nolinkurl{#1}}

\title{Excluding Parameter Three at Order Thirteen in the
\texorpdfstring{$\gamma$--$\theta$}{gamma--theta} Conjecture\\
for One-Guard Eternal Domination}
\author{Alec Kriebel}
\date{July 28, 2026}

\begin{document}

\maketitle

\begin{abstract}
In the one-guard-moves eternal domination game, attacks are made only at
unoccupied vertices and exactly one adjacent guard moves to the attacked
vertex.  The \(\gamma\)--\(\theta\) conjecture asserts that
\(\gamma(G)=\gaminf(G)\) implies \(\gamma(G)=\theta(G)\).
We prove the finite theorem that no graph on 13 vertices satisfies
\[
  \gamma(G)=\gaminf(G)=3<\theta(G).
\]
The proof fixes a maximum independent triple in an optimal eternal family
and divides the search into full-response and no-full-response branches.
The full-response branch is excluded by a 9,802-variable formula and an
independently replayed RUP refutation.  In the no-full branch, human
one-guard lemmas force two distinct response types and pure physical
representatives.  A separate 1,222-variable certificate bounds the neutral
vertices by three; a final 9,802-variable formula then covers all remaining
signature multisets.  Its deletion-free proof has 156,205 RUP additions and
zero RAT lemmas.  Independent checkers reconstruct the decisive formulas
byte for byte, verify the symmetry and coverage arguments, replay the
proofs, and check satisfiable ablations and sharp controls.  This result
does not exclude parameters four or five at order 13 and does not resolve
the universal conjecture.
\end{abstract}

\section{Introduction}\label{sec:intro}

Let \(G\) be a finite simple graph.  The domination number \(\gamma(G)\)
is the minimum size of a set whose closed neighborhood is all of
\(\V(G)\).  The clique cover number \(\theta(G)\) is the minimum number of
parts in a partition of \(\V(G)\) into cliques; equivalently,
\[
  \theta(G)=\chi(\overline G).
\tag{1.1}
\]
In the standard one-guard eternal domination game, guards occupy a
dominating set.  After an attack at an unoccupied vertex, exactly one
adjacent guard moves to the attacked vertex.  The least number of guards
that can answer every finite attack sequence is \(\gaminf(G)\).

Klostermeyer and MacGillivray raised the implication
\[
  \gamma(G)=\gaminf(G)
  \quad\Longrightarrow\quad
  \gamma(G)=\theta(G)
\tag{1.2}
\]
	\cite{KlostermeyerMacGillivray2009}.  Klostermeyer and Mynhardt later
	identified a gap in that argument and reopened the implication as an
	explicit question \cite{KlostermeyerMynhardt2015}.  Subsequent literature
	refers to their conjectural formulation as the
	\(\gamma\)--\(\theta\) conjecture
	\cite{MacGillivrayMynhardtVirgile2022}.  MacGillivray, Mynhardt, and Virgile
reported an exhaustive computation finding no counterexample through
order 11 \cite{MacGillivrayMynhardtVirgile2022}.  A separate companion
paper gives a certificate-backed exclusion at order 12.  The present paper
advances one parameter slice at the next order.

\begin{theorem}[order 13, parameter three]\label{thm:main}
There is no finite simple graph \(G\) on 13 vertices such that
\[
  \gamma(G)=\gaminf(G)=3<\theta(G).
\tag{1.3}
\]
\end{theorem}

Theorem~\ref{thm:main} is unconditional on lower-order enumeration: it
excludes exactly the displayed \((n,k)=(13,3)\) slice.  In combination with
the preceding order-12 frontier and the standard structural bounds used
there, a counterexample of order 13, if one exists, has common parameter
four or five.  Neither statement proves the conjecture for all graphs.

The proof is both mathematical and certificate-backed.  Section
\ref{sec:model} fixes the exact game and derives the equality structure.
Section \ref{sec:response} proves the response-list lemmas that cover the
no-full branch.  Sections \ref{sec:full}--\ref{sec:residual} give the three
finite certificates and their coverage arguments.  Section
\ref{sec:assemble} assembles the theorem, and Section \ref{sec:repro}
records exact hashes and reproduction instructions.

\begin{remark}[status and priority]\label{rem:status}
This manuscript and its computational artifacts were produced in an
independent, AI-assisted research campaign and have not been externally
peer reviewed.  The literature audit found no prior order-13,
parameter-three exclusion, but no worldwide priority claim is made.
\end{remark}

\section{\hspace{0.5em}The exact model and equality structure}\label{sec:model}

\begin{definition}\label{def:family}
An \emph{eternal dominating family of size \(k\)} is a nonempty family
\(\mathcal D\subseteq\binom{\V(G)}k\) such that:
\begin{enumerate}
\item every \(D\in\mathcal D\) dominates \(G\); and
\item for every \(D\in\mathcal D\) and every \(r\in\V(G)-D\), there is a
      guard \(u\in D\cap N_G(r)\) for which
      \[
        (D-\{u\})\cup\{r\}\in\mathcal D.
      \tag{2.1}
      \]
\end{enumerate}
The minimum such \(k\) is \(\gaminf(G)\).
\end{definition}

Thus no occupied-vertex attacks are imposed, and a legal response changes
exactly one guard position along exactly one edge.  We do not use an
all-guards-move, total, connected, eviction, or distance variant.

\begin{proposition}[parameter chain]\label{prop:chain}
For every nonempty graph,
\[
  \gamma(G)\leq i(G)\leq\alpha(G)\leq\gaminf(G)\leq\theta(G).
\tag{2.2}
\]
Consequently, if \(\gamma(G)=\gaminf(G)=k\), then
\[
  \gamma(G)=i(G)=\alpha(G)=\gaminf(G)=k.
\tag{2.3}
\]
\end{proposition}

\begin{proof}
Every maximal independent set dominates, and a maximum independent set is
maximal.  To see \(\alpha\leq\gaminf\), attack successively at unoccupied
vertices of a maximum independent set \(I\).  No guard already on \(I\)
can answer another attack in \(I\), so every response increases the number
of guards on \(I\).  Fewer than \(|I|\) guards eventually leave an
unoccupied vertex of \(I\) with no possible response.

Given a clique partition, place one guard in each part.  An attack is
answered by moving the guard within the attacked vertex's clique.  This
proves \(\gaminf\leq\theta\), and squeezing gives (2.3).
\end{proof}

The preceding attack argument also gives a fact used repeatedly below.

\begin{lemma}[independent states are forced]\label{lem:independent-forced}
Every independent \(k\)-set belongs to every eternal family of \(k\)-sets.
\end{lemma}

\begin{proof}
Start from any family state and attack the unoccupied vertices of the
independent \(k\)-set.  Each response increases its occupancy by one, so
closure eventually reaches the whole set.
\end{proof}

We use two elementary eternal-domination facts and one standard
small-parameter consequence in the complement.

\begin{lemma}[induced-subgraph monotonicity]\label{lem:induced}
If \(J\) is a nonempty induced subgraph of \(G\), then
\[
  \gaminf(J)\leq\gaminf(G).
\tag{2.4}
\]
\end{lemma}

\begin{proof}
Fix an eternal family in \(G\), let \(m\) be the maximum number of guards
that any family state places in \(V(J)\), and retain the states attaining
\(m\).  An attack inside \(J\) cannot be answered by moving a guard from
outside \(J\), since that would produce a family state with \(m+1\) guards
in \(J\).  Hence every response stays inside \(J\) and preserves \(m\).
Projecting the retained states to \(J\) gives a closed family of \(m\)-sets.
Every projected state dominates \(J\): an unoccupied vertex has an
adjacent responding guard in \(J\), while an occupied vertex dominates
itself.  Also \(m>0\), since an attack in nonempty \(J\) would otherwise
move a guard into \(J\).  Thus the projections form an eternal family in
\(J\).
\end{proof}

\begin{lemma}[odd antiholes]\label{lem:antihole}
For every odd \(\ell\geq5\),
\[
  \gaminf(\overline{C_\ell})=3.
\tag{2.5}
\]
\end{lemma}

\begin{proof}
The upper bound follows from
\(\theta(\overline{C_\ell})=\chi(C_\ell)=3\).
Label the underlying cycle cyclically by \(0,\ldots,\ell-1\).  A pair
fails to dominate \(\overline{C_\ell}\) exactly when its two vertices have
cyclic distance two.

Suppose an eternal two-family existed.  Its independent pair \(\{0,1\}\)
belongs to the family by Lemma~\ref{lem:independent-forced}.  If
\(\{0,d\}\) belongs to the family for odd
\(1\leq d\leq\ell-6\), attack \(d+2\).  Moving the guard at \(0\), if
legal, leaves the nondominating pair \(\{d,d+2\}\).  Closure therefore
forces the other guard to move, and \(\{0,d+2\}\) belongs to the family.
Iteration reaches \(\{0,\ell-4\}\).  Attacking \(\ell-2\) then leaves a
cyclic-distance-two pair whichever guard moves.  For \(\ell=5\), this
last attack applies directly to the initial pair.
\end{proof}

\begin{lemma}[parameter two]\label{lem:k2}
If \(\alpha(J)=\gaminf(J)=2\), then \(\theta(J)=2\).
\end{lemma}

\begin{proof}
Suppose \(\theta(J)>2\), and put \(K=\overline J\).  Then
\(\omega(K)=2<\chi(K)\), so the Strong Perfect Graph Theorem
\cite{ChudnovskyRobertsonSeymourThomas2006} supplies an induced odd hole
or odd antihole.  An odd hole in \(K\) induces an odd antihole in \(J\),
contradicting Lemmas~\ref{lem:antihole} and \ref{lem:induced}.  An odd
antihole in \(K\) induces an odd hole in \(J\).  At length at least seven
this contradicts \(\alpha(J)=2\); at length five,
self-complementarity and Lemmas~\ref{lem:antihole} and
\ref{lem:induced} again contradict \(\gaminf(J)=2\).
Hence \(\chi(K)=2\).
\end{proof}

\begin{remark}
Lemma~\ref{lem:k2} is used below only to turn an exact two-guard projection
into a bipartite complement.
\end{remark}

\section{Response lists and the no-full decomposition}\label{sec:response}

Assume for the remainder of Sections \ref{sec:response}--\ref{sec:residual}
that a graph and family satisfy
\[
  |V(G)|=13,\qquad
  \gamma(G)=\alpha(G)=\gaminf(G)=3<\theta(G).
\tag{3.1}
\]
Put \(H=\overline G\), fix an eternal family \(\F\) of triples, and fix a
maximum independent state
\[
  S=\{a,b,c\}\in\F
\tag{3.2}
\]
using Lemma~\ref{lem:independent-forced}.  For \(x\notin S\), define its
family-response list and anchor signature by
\[
  L(x)=\{u\in S:S-u+x\in\F\},
  \qquad
  \sigma(x)=N_H(x)\cap S.
\tag{3.3}
\]
An attack at \(x\) from \(S\) proves \(L(x)\ne\varnothing\).  Moreover
\[
  L(x)\subseteq S-\sigma(x),
\tag{3.4}
\]
because \(S-u+x\) must dominate the omitted anchor \(u\), while the other
two anchors do not see \(u\).  Put
\[
  Q=\{x\notin S:\sigma(x)=\varnothing\},\qquad
  A=(V(G)-S)-Q.
\tag{3.5}
\]
The vertices in \(Q\) are \(G\)-complete to \(S\) and will be called
\emph{neutral}.

\subsection{Frozen projections and response types}

For \(u\in S\), let
\[
  W_u=\{x\notin S:u\notin L(x)\}.
\tag{3.6}
\]

\begin{lemma}[restoration]\label{lem:restoration}
If \(D\in\F\) and \(u\in S-D\), then
\[
  u\in\bigcup_{x\in D-S}L(x).
\tag{3.7}
\]
\end{lemma}

\begin{proof}
Starting afresh from \(D\), attack every vertex of
\((S-D)-\{u\}\).  A guard already on \(S\) cannot answer another attack
in \(S\), because \(S\) is independent.  These attacks therefore restore
all missing anchors except \(u\) by moving guards from \(D-S\).  The
resulting family state has the form \(S-u+x\) for some \(x\in D-S\), so
\(u\in L(x)\).
\end{proof}

\begin{lemma}[frozen projection]\label{lem:frozen}
The complement-induced graph
\[
  H[(S-\{u\})\cup W_u]
\tag{3.8}
\]
is bipartite.
\end{lemma}

\begin{proof}
Let \(R=(S-\{u\})\cup W_u\), and define
\[
 \mathcal P_u=
 \{A\in\tbinom{R}{2}:\{u\}\cup A\in\F\}.
\]
This family is nonempty because \(S-\{u\}\in\mathcal P_u\).
Take \(A\in\mathcal P_u\) and attack \(r\in R-A\) from
\(\{u\}\cup A\).  If \(u\) answered, the successor \(A\cup\{r\}\)
would omit \(u\), while every one of its outside vertices belongs to
\(W_u\).  Lemma~\ref{lem:restoration} applied to that successor would
force one of those outside vertices to have \(u\) in its response list,
a contradiction.  Hence a guard in \(A\) moves, and the resulting
two-set remains in \(\mathcal P_u\).  The same response supplies an
adjacent guard for every unoccupied vertex of \(R\), so every member of
\(\mathcal P_u\) dominates \(G[R]\).  The projected states therefore form
an eternal two-family in the subgraph of \(G\) induced by \(R\).

The set \(S-\{u\}\) is independent there, so that induced graph has
\(\alpha=\gaminf=2\).  Lemma~\ref{lem:k2} gives clique cover number two;
equivalently, (3.8) is bipartite.
\end{proof}

\begin{lemma}[singleton safety]\label{lem:singleton}
If \(L(x)=L(y)=\{u\}\), then \(xy\in E(G)\).
\end{lemma}

\begin{proof}
If \(xy\in E(H)\), attack \(y\) from \(S-u+x\).  The guard at \(x\) cannot
move, so an anchor \(v\in S-\{u\}\) moves to \(y\), producing
\(S-\{u,v\}+\{x,y\}\).  Attack the now-unoccupied \(u\).  No remaining
anchor can answer.  A response by \(x\) or \(y\) produces a direct swap at
the reference state and puts \(v\) into \(L(y)\) or \(L(x)\), contrary to
both lists being \(\{u\}\).
\end{proof}

Say that \emph{type \(u\)} occurs if some vertex has the exact two-list
\[
  L(x)=S-\{u\}.
\tag{3.9}
\]

\begin{lemma}[at least two types]\label{lem:two-types}
If every outside list has size at most two, then at least two distinct
two-list types occur.
\end{lemma}

\begin{proof}
Suppose at most one type occurs, and choose \(c\in S\) so that the
occurring type, if any, is type \(c\).  Every vertex outside
\((S-\{c\})\cup W_c\) has \(c\in L(x)\).  Under the no-full hypothesis it
must have singleton list \(\{c\}\), since any two-list containing \(c\)
would have omitted color different from \(c\).  Therefore
\[
 V(H)=
 \bigl((S-\{c\})\cup W_c\bigr)
 \mathbin{\dot\cup}
 \bigl(\{c\}\cup\{x:L(x)=\{c\}\}\bigr).
\tag{3.10}
\]
The first induced part is bipartite by Lemma~\ref{lem:frozen}.  The second
is independent in \(H\): response membership gives the edges \(cx\) in
\(G\), and Lemma~\ref{lem:singleton} gives every edge between its outside
vertices.  Thus \(H\) is 3-colorable, contradicting
\(\theta(G)=\chi(H)>3\).
\end{proof}

\subsection{Physical representatives}

\begin{lemma}[two-response replication]\label{lem:replication}
Let \(x\in Q\) and suppose \(a,b\in L(x)\).  There is a vertex \(z\in A\)
such that
\[
  L(z)=\{a,b\},\qquad \sigma(z)=\{c\}.
\tag{3.11}
\]
\end{lemma}

\begin{proof}
Because \(\gamma(G)\geq3\), the pair \(\{c,x\}\) does not dominate.
Choose \(y\in N_H(c)\cap N_H(x)\).  The retained states
\(\{b,c,x\}\) and \(\{a,c,x\}\) both dominate \(y\), forcing
\(ay,by\in E(G)\), hence \(\sigma(y)=\{c\}\).

The pair \(\{c,y\}\) also fails to dominate, so choose
\(z\in N_H(c)\cap N_H(y)\).  If \(z\) missed \(a\) in \(G\), domination
of \(\{a,c,x\}\) would force \(xz\in E(G)\).  Attacking \(z\) from
\(\{b,c,x\}\) would then leave no response: \(c\) cannot move; moving
\(x\) leaves a state missing \(a\); and moving \(b\) leaves a state
missing \(y\).  The case in which \(z\) misses \(b\) is symmetric.
Therefore \(\sigma(z)=\{c\}\).

Attack \(z\) from each of \(\{b,c,x\}\) and \(\{a,c,x\}\).  The guard at
\(c\) cannot move, and moving the other anchor again leaves
\(\{c,x,z\}\), which misses \(y\).  Closure forces the guard at \(x\) to
move in both attacks.  Thus \(xz\in E(G)\) and both
\(S-a+z,S-b+z\) belong to \(\F\).  Equation (3.4) excludes \(c\) from
\(L(z)\), proving (3.11).
\end{proof}

\begin{corollary}[physical representative]\label{cor:physical}
Every occurring type \(u\) has a representative \(z_u\) satisfying
\[
  L(z_u)=S-\{u\},\qquad \sigma(z_u)=\{u\}.
\tag{3.12}
\]
\end{corollary}

\begin{proof}
Let \(t\) have list \(S-\{u\}\).  If \(ut\in E(H)\), take \(z_u=t\).
Otherwise (3.4) shows that \(t\in Q\), and
Lemma~\ref{lem:replication}, after relabeling the anchors, applies.
\end{proof}

\begin{lemma}[pure-signature doubling]\label{lem:doubling}
For every occurring type \(u\), there are distinct vertices \(z_u,w_u\)
such that
\[
 \sigma(z_u)=\sigma(w_u)=\{u\},
 \qquad z_uw_u\in E(H),
 \qquad L(z_u)=S-\{u\}.
\tag{3.13}
\]
\end{lemma}

\begin{proof}
Take \(z_u\) from Corollary~\ref{cor:physical}.  Since \(\{u,z_u\}\) does
not dominate, choose
\(w_u\in N_H(u)\cap N_H(z_u)\).  It lies outside \(S\) and has \(u\) in
its anchor signature.  If another anchor \(r\) also lay in
\(\sigma(w_u)\), let \(h\) be the third anchor.  The retained state
\(S-h+z_u=\{u,r,z_u\}\) would miss \(w_u\), a contradiction.
Thus \(\sigma(w_u)=\{u\}\), and the chosen common-neighbor incidence gives
the required \(H\)-edge.
\end{proof}

By Lemmas \ref{lem:two-types} and \ref{lem:doubling}, after permuting the
three anchors and relabeling outside vertices, every no-full candidate has
\[
\begin{array}{c|c|c}
\text{vertices}&\text{signature}&\text{additional condition}\\ \hline
3,4&\{0\}&34\in E(H),\quad L(3)=\{1,2\},\\
5,6&\{2\}&56\in E(H),\quad L(5)=\{0,1\}.
\end{array}
\tag{3.14}
\]
This is the structural input to the final certificates.

\section{The full-response branch}\label{sec:full}

A vertex \(x\notin S\) is a \emph{full target} if \(L(x)=S\).

\begin{theorem}[certified full-response exclusion]\label{thm:full}
No object satisfying (3.1) has a full target at \(S\).
\end{theorem}

\begin{proof}[Certificate and coverage proof]
Relabel \(S=012\) and a specified full target as \(x=3\).  Edge variables
encode \(H=\overline G\).  The formula has four variable families:
\[
\begin{array}{l|r}
H\text{-edges}&78\\
\text{pair common-neighbor selectors}&858\\
\text{retained triples}&286\\
\text{one-guard response selectors}&8,580\\ \hline
\text{total}&9,802.
\end{array}
\tag{4.1}
\]
Its 85,409 clauses enforce: no \(K_4\) in \(H\); an outside common
\(H\)-neighbor for every pair; domination of every retained triple; exact
one-guard closure for each unoccupied attack; the independent retained
anchor triangle and three full successors; and non-3-colorability of
\(H\).

The coloring block is complete.  A proper 3-coloring gives three distinct
colors to the anchor triangle; permute their names to \(0,1,2\).  The
formula blocks all \(3^{10}=59,049\) assignments on the remaining
vertices by requiring a monochromatic \(H\)-edge in each assignment.

Vertices \(4,\ldots,12\) are otherwise interchangeable.  Sort their
four-bit \(H\)-adjacency signatures to \(0,1,2,3\).  The eight adjacent
comparators use \(8\binom{16}{2}=960\) clauses.  Every nonsorter clause
family is covariant under the full \(S_9\) action, including the complete
coloring bank, so every unrestricted object has a sorted representative.
An independent truth table checked all 2,048 local signature pairs.

A clean-room generator reconstructed the 4,808,845-byte DIMACS file byte
for byte.  Its SHA-256 is
\begin{quote}\small
\filehash{d5a2f17ad6e61cb7ca5cb9d2930b6a0738fec32ee1d9956207dc67bb297dcb13}.
\end{quote}
The 19,874,489-byte DRAT proof has SHA-256
\begin{quote}\small
\filehash{653b01e904b97c01bfa25fbbea29fbadee603918dbaff0ea41b7ad09460fb910}.
\end{quote}
Independent replay reported \texttt{s VERIFIED} and zero RAT lemmas; a
separately retained reduced proof also replayed.  Since any hypothetical
object maps to a model of the reconstructed formula, its certified
unsatisfiability proves the theorem.
\end{proof}

Removing the 59,049 coloring clauses makes the formula satisfiable.  The
decoded graph has graph6 record
\verb!LF\|ul\XzVsaqJ! in the source
archive, parameters
\((\gamma,\alpha,\gaminf,\theta)=(3,3,3,3)\), and all 157 dominating
triples in its greatest eternal family.  This positive control shows that
the full-response and equality clauses are not intrinsically inconsistent.

\section{Four neutral vertices are impossible}\label{sec:neutral}

\begin{theorem}[four-neutral/two-port exclusion]\label{thm:four-neutral}
On 13 vertices, there is no graph with \(\gamma(G)\geq3\), an independent
retained triple \(S=\{h,i,j\}\), four distinct neutral vertices, and two
further distinct vertices \(x,y\) satisfying
\[
  \{h,j\}\subseteq L(x),\qquad
  \{h,i\}\subseteq L(y),
\tag{5.1}
\]
that admits an eternal family of triples.
\end{theorem}

\begin{proof}[Certificate and scope proof]
The formula deliberately does not impose a clique-cover gap,
connectedness, no-full lists, negative response memberships, signatures of
\(x,y\), or greatest-family status.  It uses 78 edge variables, 286 family
variables, and 858 common-neighbor selectors.  Exact one-guard closure is
expanded directly into eight clauses per retained-state/attack pair.  The
total is
\[
  1,222\text{ variables},\qquad
  24,694\text{ unique, nontautological clauses}.
\tag{5.2}
\]
An independent reconstructor produced the exact 447,906-byte instance with
SHA-256
\begin{quote}\small
\filehash{3d1a1379eb2a90ffd399e5a830b1a81881ed527c6e9db06574a390085cb5c1e0}.
\end{quote}
Its deletion-free proof contains 78,697 additions, ends in the empty
clause, and has SHA-256
\begin{quote}\small
\filehash{c4f1989ac80474a86b75ba939e494bde5928b2727fd61297eb695f3937222eee}.
\end{quote}
Strict forward replay verified every addition by reverse unit propagation,
reported zero RAT lemmas, and a fresh DRAT-to-LRAT conversion passed an
independent LRAT checker.

There is no symmetry assumption in this formula.  Any object in the
theorem can be relabeled to its named anchors, ports, and four neutral
vertices, so the certificate covers the stated universe directly.
\end{proof}

\begin{corollary}\label{cor:q3}
Every no-full candidate satisfying (3.1) has
\[
  |Q|\leq3,\qquad |A|\geq7.
\tag{5.5}
\]
\end{corollary}

\begin{proof}
Choose the two types and their distinct representatives in (3.14).  Their
two response pairs overlap in one anchor.  If four neutral vertices
existed, Theorem~\ref{thm:four-neutral} would apply after relabeling.
There are ten vertices outside \(S\), so the bound on \(A\) follows.
\end{proof}

The neutral bound is sharp with respect to the certificate hypotheses.  A
checked 13-vertex control has graph6 record
\verb!LDZZa^g|fkw[iH!,
\[
  \gamma=i=\alpha=\gaminf=\theta=3,
\tag{5.6}
\]
and a greatest eternal family of 139 triples.  It realizes three neutral
vertices and both positive port pairs.

\section{The residual no-full certificate}\label{sec:residual}

\begin{theorem}[certified residual exclusion]\label{thm:residual}
No no-full object satisfying (3.1) exists.
\end{theorem}

\begin{proof}[Formula coverage and certificate]
Use the normalization (3.14).  Label the six interchangeable remaining
vertices \(7,\ldots,12\).  Sort their three-bit anchor
signatures in nondecreasing order.  Since the four named vertices are
nonneutral, Corollary~\ref{cor:q3} leaves at most three neutral vertices
among the six residual labels.  The fourth residual vertex, label 10,
therefore has nonzero signature.

This normalization covers every candidate.  There are six ordered choices
of two distinct omitted colors, all transported to the named pair by one
of the six anchor permutations.  For the residual vertices, sorting covers
all
\[
  \binom{8+6-1}{6}=1,716
\tag{6.1}
\]
signature multisets, including ties.  The label-10 condition is equivalent
to the at-most-three-neutral bound in this sorted order.

The exact formula again has 9,802 variables.  Its 84,614 clauses have the
following independently reconstructed census:
\[
\begin{array}{l|r@{\qquad}l|r}
\text{anchor \(H\)-triangle}&3&
\text{anchor state retained}&1\\
\text{no \(H\)-\(K_4\)}&715&
\text{pair witness choice/edges}&78+1,716\\
\text{selected-state domination}&2,860&
\text{selected-state response}&2,860\\
\text{move uses a \(G\)-edge}&8,580&
\text{move reaches a retained state}&8,580\\
\text{anchored non-3-colorability}&59,049&
\text{no-full target clauses}&10\\
\text{named lists, signatures, edges}&6+12+2&
\text{residual signature sorter}&140\\
\text{at most three neutral residuals}&1&
\text{family nonempty}&1.
\end{array}
\tag{6.2}
\]
Response variables name one occupied guard and one unoccupied attack.
Their implications require the corresponding \(G\)-edge and the unique
one-guard successor.  Thus multiple true response variables represent
alternative legal moves, never simultaneous guard motion.

The clean-room reconstruction is byte-identical to the 4,784,714-byte
DIMACS file.  Its SHA-256 is
\begin{quote}\small
\filehash{76ff2768c7afd95ee535f8684515b0b15319b1f5ca69085447a1f7eba66393e1}.
\end{quote}
The retained raw DRAT trace was replayed both with and without honoring
deletion lines.  Removing those lines gives a proof containing exactly
156,205 additions, 8,878,465 bytes, and SHA-256
\begin{quote}\small
\filehash{c985ce0a602a91a0d323594e3aeecf210fa5131027ef4b6c9b6e4d4b628f1848}.
\end{quote}
Strict forward, warning-fatal replay reported \texttt{s VERIFIED} and zero
RAT lemmas.  The clean-room audit separately checked variable ranges,
clause multiplicities, absence of duplicates and tautologies, all six
anchor normalizations, all 1,716 signature multisets, and the exact
sorter truth table.  Hence no normalized candidate survives.
\end{proof}

Removing the 59,049 clique-cover clauses makes the residual formula
satisfiable, and the retained assignment satisfies every remaining
clause.  Thus the certificate does not merely prove that the equality and
structural normalization are inconsistent.

\section{Assembly and exact scope}\label{sec:assemble}

\begin{proof}[Proof of Theorem~\ref{thm:main}]
Suppose (1.3) holds.  Proposition~\ref{prop:chain} gives
\[
  \gamma=i=\alpha=\gaminf=3.
\]
Choose an optimal eternal family and a maximum independent triple \(S\).
Every outside response list is nonempty.  If some list is full,
Theorem~\ref{thm:full} is a contradiction.  Otherwise every list has size
at most two, and Theorem~\ref{thm:residual} is a contradiction.  The two
branches exhaust all possibilities.
\end{proof}

\begin{corollary}\label{cor:frontier}
Relative to the published through-order-11 result and the separately
certified order-12 theorem, a counterexample of order 13, if one exists,
has common parameter four or five.
\end{corollary}

\begin{proof}
The parameter chain and Lemma~\ref{lem:k2} give common parameter at least
three; Theorem~\ref{thm:main} excludes three.  The no-simplicial-vertex
theorem in the order-12 companion paper gives minimum degree at least two
for a minimum counterexample.  The standard domination bound for connected
minimum-degree-two graphs \cite{McCuaigShepherd1989} then gives
\(\gamma(G)\leq\lfloor2|V(G)|/5\rfloor=5\) at order 13.  Thus only four
and five remain.
\end{proof}

\begin{remark}[what has not been proved]
Finite verification does not resolve (1.2).  In particular, this paper
does not exclude order-13 parameters four or five, any parameter at order
14, or parameter three at arbitrary order.  It supplies no counterexample
and makes no claim about the all-guards-move model.
\end{remark}

\section{Reproduction and independent checks}\label{sec:repro}

The source archive is
\[
\href{https://github.com/AlecKriebel/Math}
     {\texttt{github.com/AlecKriebel/Math}},
\]
with frozen release tag
\[
\href{https://github.com/AlecKriebel/Math/releases/tag/gamma-theta-order13-k3-v1.0.0}
     {\texttt{gamma-theta-order13-k3-v1.0.0}}.
\]
From the campaign directory, the compact aggregate replay is:
\begin{quote}
\small
\texttt{python3 -I -B -W error repro/c097/replay.py}
\end{quote}
It binds the decisive hashes, replays the four-neutral and residual proofs,
reconstructs both formulas, verifies the normalization counts, checks the
sharp three-neutral equality control, and checks the satisfiable
theta-gap ablation.

The branch-specific independent checkers are:
\begin{quote}
\small
\texttt{python3 -I -B -W error reviews/order13\_full\_target\_hostile/checker.py}

\texttt{python3 -I -B -W error reviews/tight\_micro\_hostile\_review/checker.py}

\texttt{python3 -I -B -W error reviews/order13\_no\_full\_a7\_hostile/checker.py}
\end{quote}

\begin{table}[ht]
\centering
\small
\begin{tabular}{llll}
\toprule
Branch & variables & clauses & decisive proof type\\
\midrule
full response & 9,802 & 85,409 & RUP-only DRAT, plus reduced replay\\
four neutral & 1,222 & 24,694 & 78,697 addition-only RUP steps\\
residual no-full & 9,802 & 84,614 & 156,205 addition-only RUP steps\\
\bottomrule
\end{tabular}
\caption{The three finite certificates used in the proof.}
\end{table}

The discovery generators do not serve as their own decisive validators.
The review programs use separately written allocation and clause-generation
logic, inspect the graph/complement direction, truth-table the closure and
sorter gadgets, and replay pinned proof bytes with independently hashed
checkers.  Satisfiable controls are checked by direct clause evaluation and,
where graph parameters are reported, by a separate ordinary-set evaluator.

\section*{AI-assistance and verification disclosure}

The exploratory mathematics, programs, certificate design, adversarial
audits, literature review, manuscript, and publication materials were
developed with heavy assistance from ChatGPT 5.6 Sol under Alec Kriebel's
direction.  Alec Kriebel is the human author and project lead.  No outside
individual was contacted for this project, and no external expert has
reviewed the result.  No claim was accepted on model output alone: the
finite statements are bound to exact formulas, proof objects, independent
reconstructors, and direct controls.  Those checks are evidence about the
encoded theorem; they are not peer review.

\section*{Acknowledgment of prior work}

The through-order-11 computation is due to Gary MacGillivray, C. M.
Mynhardt, and V. Virgile.  The original question is due to William F.
Klostermeyer and Gary MacGillivray, and its later formulation as the
\(\gamma\)--\(\theta\) conjecture is due to William F. Klostermeyer and
C. M. Mynhardt.

\bibliographystyle{plain}
\bibliography{references}

\end{document}
